ShAI is a procedural, CLI-driven application that combines local LLM access, FAISS-based retrieval,
and Ollama container management. The architecture is organized by package and file, with each module
owning a small, explicit responsibility.
\label{a-shai-architecture}

\subsection*{Package tree}

\begin{verbatim}
shAI_ostap4ello/src/
  __main__.py
  config/
  db/
  interpreter/
  llm/
  utils/
  workflows/
\end{verbatim}


\subsection*{Component summaries}

Modules in the applications can be divided into two layers: \textsl{low-level}
and \textsl{high-level}. The low-level layer modules are mostly independent,
providing single-responsibility functions, while high-level modules reuse this
functional and operate on it, providing more complex workflows.

\texttt{src/<module>/\_\_init\_\_.py} Re-exports the public functions provided
by the \texttt{module}.


\subsubsection*{High-level modules}

\paragraph{\texttt{src/\_\_main\_\_.py}}
Main CLI entrypoint. Loads config, resolves defaults, dispatches subcommands.
Reuses functions from \texttt{src/<module>/\_\_main\_\_.py}. Mainly calls functions from
src/workflows or low-level modules directly.

\paragraph{\texttt{src/config/}}
Implements centralized configuration loading and default propagation. It reads
\texttt{\textasciitilde/.config/shai/config.conf}, resolves defaults such as the database path and
Ollama URL, and updates module-level values so the rest of the application sees one consistent
runtime configuration.

\paragraph{\texttt{src/workflows/}}
This module contains high-level workflows, which reuse low-level modules.

\texttt{src/workflows/rag.py}:
Builds the retrieval-augmented generation flow. It converts search results into readable context,
loads referenced file content, and asks the LLM to answer using the retrieved documents.

\texttt{src/workflows/classifier\_bash.py}:
Implements the Bash-vs-natural-language classifier.


\subsubsection*{Low-level modules}

\paragraph{\texttt{src/llm/}}
Contains the core OpenAI-compatible Ollama client helpers. It provides client construction,
non-streaming generation, streaming generation, single-string embeddings, and batched embeddings.

\paragraph{\texttt{src/db/\_\_main\_\_.py}}
Implements the CLI for document indexing and retrieval. It wires the build, search, and check
commands, resolves path options, and connects CLI flags to the database layer.

\paragraph{\texttt{src/db/}}
Provides database capabilities: it builds FAISS indexes, performs search, exposes metadata
inspection.

\texttt{src/db/utils/}
Contains helper utilities for database operation.

\paragraph{\texttt{src/utils/}}
The module provides various useful functions, not related to application
workflows.

\paragraph{\texttt{src/utils/adapter.py}}
Provides the system integration layer to run various Bash scripts. It runs the
Ollama Docker helper script, starts and stops the container, checks container
status, waits for API readiness, and manages models through the Ollama HTTP API.

\paragraph{\texttt{src/*.py}}
Contains various functions exported by the module.

\paragraph{\texttt{src/*.bash} or \texttt{src/*.sh}}
A collection of shell scripts called and exported by the
\texttt{src/utils/adapter.py}.


\subsection*{Runtime flow}

\begin{enumerate}
    \item The user invokes \texttt{shai} or a module-specific command.
    \item The relevant \texttt{\_\_main\_\_.py} parser resolves arguments and configuration defaults.
    \item Database commands index or search documents through the FAISS-backed storage layer.
    \item Workflow commands combine retrieval and generation to produce a final answer.
    \item Utility commands manage Ollama and related local tooling.
    \item The shared configuration layer keeps the paths, models, and service endpoints aligned.
\end{enumerate}


\subsection*{Design characteristics}

\begin{itemize}
    \item \textbf{Procedural style}: logic is organized as functions and modules rather than classes.
    \item \textbf{Local-first execution}: all model access goes through a local Ollama service.
    \item \textbf{Explicit layering}: each package owns one part of the stack, from CLI parsing to
    retrieval to generation.
    \item \textbf{Shared defaults}: configuration is centralized so every command sees the same runtime
    state.
\end{itemize}
